\documentclass{article}
\usepackage[utf8]{inputenc}
\usepackage{xcolor}
\usepackage{amsmath}
\usepackage{graphicx}
\usepackage{float}
\usepackage{minted}
\usepackage{circuitikz}
\usepackage{tikz}
\usetikzlibrary{shapes, arrows, calc, positioning, circuits.logic.US, circuits.logic.IEC}
\usepackage{geometry}
\geometry{a4paper, margin=1in}
\usepackage{hyperref}
\usemintedstyle{trac}

\usepackage{fancyhdr}
\pagestyle{fancy}
\fancyhf{}
\renewcommand{\headrulewidth}{0pt}
\fancyfoot[C]{\small 010153101 Digital Circuit and Microprocessor Fundamental}
\fancyfoot[R]{\thepage}

% Redefine plain style to match fancy style (for the title page)
\fancypagestyle{plain}{%
  \fancyhf{}%
  \fancyfoot[C]{\small 010153101 Digital Circuit and Microprocessor Fundamental}%
  \fancyfoot[R]{\thepage}%
  \renewcommand{\headrulewidth}{0pt}%
}

\title{Laboratory Exercise 5 \\ Multiplexers, Comparators and a Simple ALU}
\author{}
\date{}

\begin{document}

\maketitle

\section*{Introduction}
This laboratory exercise accompanies the second half of Lecture 5 (Combinational Circuit Building Blocks): multiplexers, Boolean function implementation with a multiplexer, tri-state gates, the Arithmetic Logic Unit (ALU) (and, as an extension, the magnitude comparator). The circuits are written in VHDL. You first build a 2-to-1 multiplexer from basic gates on the breadboard, the same circuit that you will then describe in VHDL. You will build a 4-to-1 multiplexer from 2-to-1 multiplexers, check it with a self-checking \textbf{testbench}, and then combine the building blocks of Labs 3 to 5 into a small ALU. The ALU is described with a first VHDL \texttt{PROCESS}, which prepares you for the sequential circuits in Lab 6.

\section*{Objectives}
\begin{enumerate}
    \item To build and test a 2-to-1 multiplexer from basic gates (NOT, AND, OR) on the breadboard.
    \item To build a larger multiplexer from smaller multiplexers by structural VHDL, and to compare it with a behavioral description.
    \item To verify a design exhaustively with a self-checking VHDL testbench.
    \item To implement a Boolean function with a multiplexer.
    \item (Extension) To design a 4-bit magnitude comparator.
    \item To describe a combinational circuit with a \texttt{PROCESS} and a \texttt{CASE} statement, and to use a tri-state output.
\end{enumerate}

\section*{Equipment Required}
\begin{itemize}
    \item Breadboard, jumper wires, $+5\,\mathrm{V}$ supply, DIP switches, LEDs, resistors ($3\times 330\,\Omega$ for the switches and one for the LED), multimeter and oscilloscope (optional)
    \item TTL ICs: 74LS04 (hex inverter), 74LS08 (quad AND), 74LS32 (quad OR)
    \item Intel DE0-CV board and a PC with Quartus Prime (and the simulator selected by your instructor)
    \item Files from previous labs: \texttt{hex7seg.vhd} (Lab 4)
\end{itemize}

\section*{Pre-lab (complete before the session)}
The session is two hours long, so complete the following \emph{before} the session.
\begin{enumerate}
    \item Complete Table \ref{tab:mux_truth} (expected $Y$ of a 2-to-1 multiplexer, $Y = \bar{S}\,I_0 + S\,I_1$), and check the pin numbers of the 74LS04, 74LS08 and 74LS32 in Table \ref{tab:mux_wiring}.
    \item Draw a 4-to-1 multiplexer built with three 2-to-1 multiplexers. Which select signal drives the first stage and which the second stage?
    \item For $F(A,B,C) = \sum m(1,3,5,6,7)$ fill in the inputs of a 4-to-1 multiplexer with $A$, $B$ on the select inputs ($S_1 = A$, $S_0 = B$): $I_0 =$ \underline{\hspace{1cm}}, $I_1 = $ \underline{\hspace{1cm}}, $I_2 = $ \underline{\hspace{1cm}}, $I_3 = $ \underline{\hspace{1cm}}.
\end{enumerate}

\section*{Part I: A 2-to-1 Multiplexer on the Breadboard (about 20 min)}

A 2-to-1 multiplexer has two data inputs ($I_0$, $I_1$), a select input $S$ and an output $Y$. If $S=0$ then $Y=I_0$, otherwise $Y=I_1$:
\[ Y = \bar{S}\cdot I_0 + S\cdot I_1 \]
Build it with the gates of Figure \ref{fig:mux_logic}.

\begin{figure}[H]
    \centering
    \begin{circuitikz}[circuitikz/logic ports=ieee]
        \draw
        (0,2) node[and port] (and1) {}
        (0,0) node[and port] (and2) {}
        (2.5,1) node[or port] (or1) {};

        % NOT gate output aligned with AND1 input 2
        \draw (and1.in 2) -- ++(-0.5,0) node[not port, anchor=out, scale=0.5] (not1) {};

        % I0 on AND1 input 1, I1 on AND2 input 2
        \draw (and1.in 1) -- ++(-2.5,0) node[label=left:$I_0$] (I0) {};
        \draw (and2.in 2) -- ++(-2.5,0) node[label=left:$I_1$] (I1) {};

        % S junction
        \draw (not1.in) -- ++(-0.5,0) coordinate (s_split) node[circ] {};
        \draw (s_split) -- (s_split -| I0) node[label=left:$S$] (S) {};
        \draw (s_split) |- (and2.in 1);

        % AND outputs to OR inputs
        \draw (and1.out) -| (or1.in 1);
        \draw (and2.out) -| (or1.in 2);

        % Output Y
        \draw (or1.out) -- (3.5,1) node[right] {$Y$};
    \end{circuitikz}
    \caption{Logic diagram of the 2-to-1 multiplexer.}
    \label{fig:mux_logic}
\end{figure}

\begin{enumerate}
    \item Build the three switch circuits for $S$, $I_0$, $I_1$ as in Lab 0 (switch to $+5\,\mathrm{V}$, $330\,\Omega$ pull-down to GND). Note that $S$ drives two gate inputs ($2\times 0.4\,\mathrm{mA}$ through $330\,\Omega$ is about $0.26\,\mathrm{V}$, below $V_{IL}$, so the pull-down is still fine).
    \item Wire the circuit according to Table \ref{tab:mux_wiring}. Connect pin 14 to $+5\,\mathrm{V}$ and pin 7 to GND on every IC. The output $Y$ goes to the LED indicator stage of Lab 0 (an inverter of the 74LS04 with the LED and $330\,\Omega$ to $+5\,\mathrm{V}$), so the LED is ON when $Y=1$.
    \item Set the switches to all eight combinations and complete Table \ref{tab:mux_truth}.
\end{enumerate}

\begin{table}[H]
    \centering
    \begin{tabular}{|l|l|}
        \hline
        \textbf{From} & \textbf{To} \\
        \hline
        Switch $S$ & 74LS04 pin 1 and 74LS08 pin 4 \\ \hline
        74LS04 pin 2 ($\bar{S}$) & 74LS08 pin 2 \\ \hline
        Switch $I_0$ & 74LS08 pin 1 \\ \hline
        Switch $I_1$ & 74LS08 pin 5 \\ \hline
        74LS08 pin 3 ($\bar{S} I_0$) & 74LS32 pin 1 \\ \hline
        74LS08 pin 6 ($S I_1$) & 74LS32 pin 2 \\ \hline
        74LS32 pin 3 ($Y$) & 74LS04 pin 5 (LED indicator stage) \\ \hline
        74LS04 pin 6 & LED cathode (anode through $330\,\Omega$ to $+5\,\mathrm{V}$) \\ \hline
    \end{tabular}
    \caption{Wiring of the 2-to-1 multiplexer (74LS04, 74LS08 and 74LS32).}
    \label{tab:mux_wiring}
\end{table}

\begin{table}[H]
    \centering
    \begin{tabular}{|c|c|c|c|c|}
        \hline
        $S$ & $I_1$ & $I_0$ & $Y$ expected & LED observed \\
        \hline
        0&0&0&&\\ \hline
        0&0&1&&\\ \hline
        0&1&0&&\\ \hline
        0&1&1&&\\ \hline
        1&0&0&&\\ \hline
        1&0&1&&\\ \hline
        1&1&0&&\\ \hline
        1&1&1&&\\ \hline
    \end{tabular}
    \caption{Function table of the 2-to-1 multiplexer.}
    \label{tab:mux_truth}
\end{table}

\section*{Part II: Multiplexers in VHDL and a Testbench (about 30 min)}

\subsection*{2-to-1 multiplexer}
The 2-to-1 multiplexer with $Y = \bar{S}\,I_0 + S\,I_1$ is the circuit that you have just built. In VHDL, the \texttt{WHEN ... ELSE} conditional assignment describes the same behavior in one line.

\begin{figure}[H]
    \centering
    \begin{minted}[frame=single, fontsize=\small, linenos]{vhdl}
LIBRARY ieee;
USE ieee.std_logic_1164.ALL;

ENTITY mux2to1 IS
    PORT ( I0, I1, S : IN  STD_LOGIC;
           Y         : OUT STD_LOGIC );
END mux2to1;

ARCHITECTURE behavior OF mux2to1 IS
BEGIN
    Y <= I0 WHEN S = '0' ELSE I1;
END behavior;
    \end{minted}
    \caption{A 2-to-1 multiplexer.}
    \label{fig:mux2}
\end{figure}

\subsection*{4-to-1 multiplexer in two ways}
\begin{enumerate}
    \item Create a new Quartus project \texttt{lab5} and enter \texttt{mux2to1.vhd} (Figure \ref{fig:mux2}).
    \item Write \texttt{mux4\_struct}: a 4-to-1 multiplexer with inputs $I(3..0)$, $S(1..0)$ and output $Y$ using \textbf{three instances} of \texttt{mux2to1}. Use $S_0$ in the first stage and $S_1$ in the second stage (see Pre-lab). You need two internal signals for the outputs of the first stage.
    \item Write \texttt{mux4\_behav}: the same multiplexer described with a selected signal assignment:
\begin{minted}[frame=single, fontsize=\small]{vhdl}
WITH S SELECT
    Y <= I(0) WHEN "00",
         I(1) WHEN "01",
         I(2) WHEN "10",
         I(3) WHEN OTHERS;
\end{minted}
\end{enumerate}

\subsection*{A self-checking testbench}
A testbench is a VHDL entity without ports that instantiates the design under test, applies inputs, and checks the outputs. The testbench in Figure \ref{fig:tb} applies all $2^6 = 64$ combinations of $I$ and $S$ to both multiplexers and reports an error if the outputs differ. It uses a \texttt{PROCESS}, a statement whose body is executed \emph{in sequence}, with \texttt{WAIT FOR} to let time pass.

\begin{figure}[H]
    \centering
    \begin{minted}[frame=single, fontsize=\small, linenos]{vhdl}
LIBRARY ieee;
USE ieee.std_logic_1164.ALL;
USE ieee.numeric_std.ALL;

ENTITY mux4_tb IS END mux4_tb;

ARCHITECTURE sim OF mux4_tb IS
    SIGNAL I          : STD_LOGIC_VECTOR(3 DOWNTO 0);
    SIGNAL S          : STD_LOGIC_VECTOR(1 DOWNTO 0);
    SIGNAL Y_s, Y_b   : STD_LOGIC;
BEGIN
    u1 : ENTITY work.mux4_struct PORT MAP (I => I, S => S, Y => Y_s);
    u2 : ENTITY work.mux4_behav  PORT MAP (I => I, S => S, Y => Y_b);

    stimulus : PROCESS
        VARIABLE v : UNSIGNED(5 DOWNTO 0);
    BEGIN
        FOR k IN 0 TO 63 LOOP
            v := TO_UNSIGNED(k, 6);
            I <= STD_LOGIC_VECTOR(v(5 DOWNTO 2));
            S <= STD_LOGIC_VECTOR(v(1 DOWNTO 0));
            WAIT FOR 10 ns;
            ASSERT Y_s = Y_b REPORT "Mismatch at k = " & INTEGER'IMAGE(k)
                SEVERITY ERROR;
        END LOOP;
        REPORT "Test finished";
        WAIT;
    END PROCESS;
END sim;
    \end{minted}
    \caption{Exhaustive self-checking testbench for the two multiplexers.}
    \label{fig:tb}
\end{figure}

\begin{enumerate}
    \setcounter{enumi}{3}
    \item Compile both designs and run the simulation of \texttt{mux4\_tb} (the simulator and its set-up are explained by your instructor) for at least $640\,\mathrm{ns}$. Are there any error messages? Show the waveform of $I$, $S$, $Y\_s$ and $Y\_b$.
\end{enumerate}

\section*{Part III: Implementing a Function with a Multiplexer (about 15 min)}

A $2^n$-to-1 multiplexer can implement any function of $n+1$ variables (Shannon's expansion, Lecture 5). Use your \texttt{mux4\_behav} to implement $F(A,B,C) = \sum m(1,3,5,6,7)$.

\begin{enumerate}
    \item Write an entity \texttt{func\_mux} with inputs $A$, $B$, $C$ and output $F$. Connect $A$ to $S_1$, $B$ to $S_0$ and the inputs $I_3..I_0$ as you found in the Pre-lab (the constants are \texttt{'0'} and \texttt{'1'}, and signal $C$ for the others).
    \item Simulate \texttt{func\_mux} (or add it to a testbench like Figure \ref{fig:tb}) and check all 8 input combinations.
    \item Implement \texttt{func\_mux} on the board with $A$, $B$, $C$ = SW2--SW0 and $F$ = LEDR0 and complete Table \ref{tab:func}.
\end{enumerate}

\begin{table}[H]
    \centering
    \begin{tabular}{|c|c|c|c|c|}
        \hline
        $A$ & $B$ & $C$ & $F$ expected & $F$ measured \\
        \hline
        0&0&0&0& \\ \hline
        0&0&1&1& \\ \hline
        0&1&0&0& \\ \hline
        0&1&1&1& \\ \hline
        1&0&0&0& \\ \hline
        1&0&1&1& \\ \hline
        1&1&0&1& \\ \hline
        1&1&1&1& \\ \hline
    \end{tabular}
    \caption{Function implemented with a multiplexer.}
    \label{tab:func}
\end{table}

\section*{Part IV: A Simple ALU with Tri-State Output (about 35 min)}

An ALU performs several arithmetic or logic operations on two operands; a multiplexer selects which result appears at the output. Design a 4-bit ALU with the operations of Table \ref{tab:alu} (Figure \ref{fig:alu}). Use a \texttt{PROCESS} with a \texttt{CASE} statement for the multiplexing. In a process, signals read inside the process must appear in its \textbf{sensitivity list} (the list in parentheses after \texttt{PROCESS}).

\begin{table}[H]
    \centering
    \begin{tabular}{|c|c|}
        \hline
        \textbf{Op (2 bits)} & \textbf{Function} \\
        \hline
        00 & $Y = A + B$ \\ \hline
        01 & $Y = A - B$ \\ \hline
        10 & $Y = A$ AND $B$ \\ \hline
        11 & $Y = A$ OR $B$ \\ \hline
    \end{tabular}
    \caption{ALU operations.}
    \label{tab:alu}
\end{table}

\begin{figure}[H]
    \centering
    \begin{tikzpicture}[font=\small]
        \node[draw, minimum width=3.6cm, minimum height=2.6cm, align=center] (alu) at (0,0) {4-bit ALU\\$Y_{int}$};
        \node[draw, minimum width=1.5cm, minimum height=1.0cm, align=center] (tri) at (4.4,0) {tri-state\\buffer};
        \draw[->] (-3.5,0.7) node[left] {$A$ (4)} -- (alu.west |- 0,0.7);
        \draw[->] (-3.5,-0.7) node[left] {$B$ (4)} -- (alu.west |- 0,-0.7);
        \draw[->] (0,-2.4) node[below] {Op (2)} -- (alu.south);
        \draw[->] (alu.east) -- node[above] {$Y_{int}$} (tri.west);
        \draw[->] (tri.east) -- ++(1.2,0) node[right] {$Y$ (LEDR)};
        \draw[->] (4.4,-2.0) node[below] {$OE$} -- (tri.south);
        \draw[->] (0.9,1.3) -- ++(0,0.8) node[above] {Zero flag};
        \draw[->] (2.2,0) |- (2.2,-1.6) -- ++(0.8,0) node[right] {$Y_{int}$ to HEX0};
        \fill (2.2,0) circle (1.5pt);
    \end{tikzpicture}
    \caption{Block diagram of the ALU with a tri-state output stage.}
    \label{fig:alu}
\end{figure}

\begin{figure}[H]
    \centering
    \begin{minted}[frame=single, fontsize=\small, linenos]{vhdl}
LIBRARY ieee;
USE ieee.std_logic_1164.ALL;
USE ieee.numeric_std.ALL;

ENTITY alu4 IS
    PORT ( A, B  : IN  STD_LOGIC_VECTOR(3 DOWNTO 0);
           Op    : IN  STD_LOGIC_VECTOR(1 DOWNTO 0);
           OE    : IN  STD_LOGIC;
           Yint  : OUT STD_LOGIC_VECTOR(3 DOWNTO 0);   -- always driven
           Y     : OUT STD_LOGIC_VECTOR(3 DOWNTO 0);   -- tri-state output
           Zero  : OUT STD_LOGIC );
END alu4;

ARCHITECTURE behavior OF alu4 IS
    SIGNAL R : STD_LOGIC_VECTOR(3 DOWNTO 0);
BEGIN
    PROCESS (A, B, Op)
    BEGIN
        CASE Op IS
            WHEN "00"   => R <= STD_LOGIC_VECTOR(UNSIGNED(A) + UNSIGNED(B));
            WHEN "01"   => R <= STD_LOGIC_VECTOR(UNSIGNED(A) - UNSIGNED(B));
            WHEN "10"   => R <= A AND B;
            WHEN OTHERS => R <= A OR B;
        END CASE;
    END PROCESS;

    Yint <= R;
    Zero <= '1' WHEN R = "0000" ELSE '0';
    Y    <= R WHEN OE = '1' ELSE "ZZZZ";     -- tri-state buffer
END behavior;
    \end{minted}
    \caption{A 4-bit ALU with zero flag and tri-state output.}
    \label{fig:alu_vhdl}
\end{figure}

\begin{enumerate}
    \item Add \texttt{alu4.vhd} to the project and compile it. Why can the output ports \texttt{Yint} and \texttt{Y} not be read inside the architecture, and why do we use the internal signal \texttt{R}?
    \item Simulate the ALU. Choose test vectors that include the following cases and record the results in Table \ref{tab:alu_test}: $9+8$ (the result does not fit in 4 bits), $3-5$, and $A-A$ (the Zero flag).
    \item Write the top-level entity: $A$ = SW3--SW0, $B$ = SW7--SW4, $Op$ = SW9--SW8, $Y$ = LEDR3--LEDR0, \texttt{Zero} = LEDR9, and \texttt{Yint} displayed on \texttt{HEX0} with \texttt{hex7seg}. Use push-button KEY0 as the output enable: \texttt{OE <= NOT KEY(0);} (KEY is active-low, so the outputs are enabled \emph{while the button is pressed}).
    \item Download to the board and test Table \ref{tab:alu_test}. What do the LEDs display when KEY0 is not pressed? Why is this output state called ``high impedance''?
    \item Tri-state outputs allow several circuits to share one wire (a bus), provided that only one is enabled at any time. \emph{Inside} an FPGA chip no tri-state buses exist; the synthesis tool converts the tri-state buffers into multiplexers, except at the chip pins. Why can two ordinary gate outputs not be connected together without tri-state control? (Lecture 5, Exercise 5-3.)
\end{enumerate}

\begin{table}[H]
    \centering
    \begin{tabular}{|c|c|c|c|c|c|c|}
        \hline
        Op & $A$ & $B$ & $Y_{int}$ expected & HEX0 & Zero & LEDR3--0 (KEY0 pressed) \\
        \hline
        00 & 9 & 8 & & & & \\ \hline
        01 & 3 & 5 & & & & \\ \hline
        01 & 7 & 7 & & & & \\ \hline
        10 & 1100 & 1010 & & & & \\ \hline
        11 & 1100 & 1010 & & & & \\ \hline
    \end{tabular}
    \caption{ALU test results.}
    \label{tab:alu_test}
\end{table}

\section*{Extension (take-home, or if time permits)}
\begin{enumerate}
    \item (Part I) Measure $Y$ with the multimeter for $Y=0$ and $Y=1$ (valid TTL levels?). With $I_0=0$, $I_1=1$, connect Channel 1 of the oscilloscope to $S$ and Channel 2 to $Y$, toggle $S$ and estimate the delay of the three gate levels. Repeat with $I_0=1$, $I_1=0$.
    \item (Part II) Compile \texttt{mux4\_struct} and \texttt{mux4\_behav} separately (set each as the top-level entity) and compare the number of logic elements and the RTL Viewer pictures. Connect \texttt{mux4\_behav} to the board ($I$ = SW3--SW0, $S$ = SW9--SW8, $Y$ = LEDR0) and test all select values.
    \item (Part II) To see the effect of a failing test, change one line of \texttt{mux4\_struct} deliberately (for example, swap $S_0$ and $S_1$) and rerun the testbench. What message appears? Restore the correct code.
    \item (Part III) Write a second entity \texttt{func\_sop} that implements $F$ as a minimized SOP expression found with a K-map ($F = $ \underline{\hspace{4cm}}), show that it is equivalent to \texttt{func\_mux}, and compare the number of logic elements.
\end{enumerate}

\subsection*{4-bit Magnitude Comparator}

A magnitude comparator compares two numbers $A$ and $B$. For 4 bits, define the \emph{equality} signals $x_i = A_i \odot B_i$ (XNOR). Then
\begin{align*}
(A = B) &= x_3\,x_2\,x_1\,x_0 \\
(A > B) &= A_3\bar{B}_3 + x_3 A_2\bar{B}_2 + x_3x_2A_1\bar{B}_1 + x_3x_2x_1A_0\bar{B}_0 \\
(A < B) &= \bar{A}_3 B_3 + x_3 \bar{A}_2 B_2 + x_3x_2\bar{A}_1 B_1 + x_3x_2x_1\bar{A}_0 B_0
\end{align*}

\begin{enumerate}
    \item Write an entity \texttt{comp4\_eq} that uses these equations with one \texttt{SIGNAL x : STD\_LOGIC\_VECTOR(3 DOWNTO 0)} for the equality signals, e.g. \texttt{x <= A XNOR B;}. The ports are $A$, $B$ (4 bits) and \texttt{AGB}, \texttt{AEB}, \texttt{ALB}.
    \item Write a second entity \texttt{comp4\_beh} with the same ports, but using the relational operators on unsigned numbers:
\begin{minted}[frame=single, fontsize=\small]{vhdl}
USE ieee.numeric_std.ALL;
...
AGB <= '1' WHEN UNSIGNED(A) >  UNSIGNED(B) ELSE '0';
\end{minted}
    \item Write a testbench (modify Figure \ref{fig:tb}: loop over $0..255$, split $v$ into $A$ and $B$) that compares both comparators for all 256 input combinations.
    \item On the board: $A$ = SW3--SW0, $B$ = SW7--SW4, \texttt{AGB}, \texttt{AEB}, \texttt{ALB} on LEDR2, LEDR1, LEDR0. Also display $A$ on \texttt{HEX1} and $B$ on \texttt{HEX0} (with \texttt{hex7seg}) and check the three cases of Table \ref{tab:comp}.
\end{enumerate}

\begin{table}[H]
    \centering
    \begin{tabular}{|c|c|c|c|c|}
        \hline
        $A$ & $B$ & AGB & AEB & ALB \\
        \hline
        1011 & 1001 & & & \\ \hline
        0110 & 0110 & & & \\ \hline
        0011 & 1000 & & & \\ \hline
    \end{tabular}
    \caption{Comparator results.}
    \label{tab:comp}
\end{table}


\section*{Check-off}
Show your instructor or teaching assistant:
\begin{enumerate}
    \item The 2-to-1 multiplexer on the breadboard (Table \ref{tab:mux_truth}).
    \item The testbench waveform of the multiplexers without errors.
    \item The function implemented with a multiplexer on the board.
    \item The ALU with tri-state output working on the board (Table \ref{tab:alu_test}).
\end{enumerate}

\vspace{1cm}
\noindent \textbf{Updated By:} R. Sutthaweekul \\
\textbf{Version:} 1.0 \\
\textbf{Release Date:} 2026-11-11

\end{document}
